\documentclass[a4paper,12pt]{extarticle}
\usepackage{geometry}
\usepackage[T1]{fontenc}
\usepackage[utf8]{inputenc}
\usepackage[english,russian]{babel}
\usepackage{amsmath}
\usepackage{amsthm}
\usepackage{amssymb}
\usepackage{fancyhdr}
\usepackage{setspace}
\usepackage{graphicx}
\usepackage{colortbl}
\usepackage{tikz}
\usepackage{pgf}
\usepackage{subcaption}
\usepackage{listings}
\usepackage{indentfirst}
\usepackage[
backend=biber,
style=numeric,
maxbibnames=99
]{biblatex}
\addbibresource{refs.bib}
\usepackage[colorlinks,citecolor=blue,linkcolor=blue,bookmarks=false,hypertexnames=true, urlcolor=blue]{hyperref}
\usepackage{indentfirst}
\usepackage{mathtools}
\usepackage{booktabs}
\usepackage[flushleft]{threeparttable}
\usepackage{tablefootnote}

\usepackage{chngcntr}
\counterwithin{table}{section}
\counterwithin{figure}{section}

\graphicspath{{graphics/}}

\geometry{left=2.5cm,right=1.0cm,top=2.0cm,bottom=2.0cm}
\setlength{\parindent}{1.25cm}
\renewcommand{\baselinestretch}{1.5}

\begin{document}
\input{title_kr}
\newpage
\setcounter{page}{2}

{
\hypersetup{linkcolor=black}
\tableofcontents
}

\newpage

\section*{Аннотация}
В данной работе рассматривается разработка чат-бота с применением технологий искусственного интеллекта для абитуриентов <<Национального исследовательского университета Высшая школа экономики>>, предназначенного для предоставления справочной информации по основным аспектам процесса поступления.

Разработка чат-бота основана на подходе Retrieval-Augmented Generation (RAG). Сформирована база знаний из официальных документов и страниц сайта НИУ ВШЭ, выполнена их предварительная обработка, разбиение на смысловые фрагменты (чанки) и представление в векторном виде для семантического поиска. Генерация ответов выполняется большими языковыми моделями (LLM) с использованием извлечённого контекста.

Результатом работы является программный прототип чат-бота, способного отвечать на типовые вопросы абитуриентов; проведены эксперименты по подбору ключевых параметров RAG-пайплайна и оценено качество системы на запечатанном тестовом наборе.

\section*{Abstract}
This paper presents the development of an AI-based chatbot for applicants of the National Research University Higher School of Economics, intended to provide reference information on the main aspects of the admission process.

The chatbot is developed using the Retrieval-Augmented Generation (RAG) approach. A knowledge base is built from official documents and web pages of HSE, preprocessed, chunked into semantic fragments, and represented as vector embeddings for relevant retrieval. Answer generation is performed by large language models (LLM) conditioned on retrieved context.

The result is a software prototype of an admission chatbot capable of answering typical applicants' questions; experiments select key RAG parameters, and quality is evaluated on a held-out test set.

\addcontentsline{toc}{section}{Аннотация}

\section*{Ключевые слова}
чат-бот, Retrieval-Augmented Generation, абитуриенты, НИУ ВШЭ, векторные базы данных

\pagebreak

\section{Введение}
Поступление в высшее учебное заведение требует от абитуриентов работы с большим количеством нормативных документов, справочных материалов, включая правила приёма, описание образовательных программ, условия конкурсного отбора, минимальные баллы и сведения о льготах. Поиск ответа на конкретный вопрос занимает значительное время и требует анализа нескольких источников.

Генеративные чат-боты на основе больших языковых моделей (LLM) способны понимать сложные вопросы и формировать связные ответы. Однако при использовании без контекста и привязки к базе знаний они могут галлюцинировать. В задаче помощи абитуриентам, где корректность и соответствие официальным правилам имеют принципиальное значение, это ограничение существенно снижает применимость LLM.

В данной работе рассматривается разработка прототипа чат-бота для абитуриентов <<Национального исследовательского университета Высшая школа экономики>>, предназначенного для предоставления справочной информации на основе официальных источников. Используется архитектура Retrieval-Augmented Generation (RAG)~\cite{ctan,rag}: запрос пользователя направляется в семантический поиск по базе знаний, найденный контекст передаётся языковой модели и используется для формирования ответа.

\section{Цели и задачи}

\subsection{Цели}
Целью работы является разработка чат-бота для информационной поддержки абитуриентов НИУ ВШЭ, обрабатывающего запросы на русском языке и формирующего ответы на основе официальных источников.

\subsection{Задачи}
\begin{itemize}
    \item Сбор и обработка данных из официальных источников.
    \item Разработка RAG-пайплайна.
    \item Подбор оптимальных параметров через серию экспериментов.
    \item Создание Telegram-интерфейса.
    \item Оценка качества на тестовом наборе.
\end{itemize}

\section{Обзор литературы}

\subsection{Подходы}

\subsubsection{FAQ}
Традиционный подход к автоматизированному консультированию --- FAQ-разделы с фиксированными парами вопрос–ответ. Они предсказуемы, но плохо масштабируются на переформулировки и требуют ручной актуализации.

\subsubsection{Генеративные LLM}
С появлением больших языковых моделей~\cite{gpt} стало возможно генерировать связные ответы на естественном языке. Без контекста LLM опираются на параметрическую память, не гарантируя актуальности и корректности в нормативных формулировках.

\subsubsection{Retrieval-Augmented Generation}
RAG~\cite{ctan,rag} объединяет семантический поиск по внешней базе документов с генеративной моделью. Современные реализации используют плотные эмбеддинги~\cite{word2vec,sbert,e5,bgem3} и специализированные векторные хранилища~\cite{chromadb}.

\subsection{Итог}
Архитектура RAG выбрана как наилучшая для задачи: обеспечивает актуальность ответов через обновляемую базу знаний и контроль источников при сохранении гибкости естественного диалога.

\section{Архитектура и обоснование технологических решений}

\subsection{Общая схема системы}
Система состоит из двух циклов: offline-конвейер (скрейпинг, очистка, чанкирование, эмбеддинг, upsert в индекс) и online-цикл (Telegram-бот $\rightarrow$ RAG-пайплайн $\rightarrow$ генератор). Компоненты: aiogram-бот~\cite{aiogram}, ChromaDB~\cite{chromadb} (embedded), эмбеддер \texttt{BAAI/bge-m3}~\cite{bgem3}, LLM через Ollama~\cite{ollama}, ingestion-скрипты на httpx, trafilatura~\cite{trafilatura}, BeautifulSoup, PyMuPDF~\cite{pymupdf}. Бот вызывает \texttt{rag.pipeline.answer(query)} напрямую как Python-модуль. Источники приёмной комиссии меняются редко, поэтому периодический пересбор индекса (порядка минут на MacBook M4) оправдан.

\subsection{Источники данных}
Корпус включает шесть типов источников: правила приёма и приложения (PDF, PyMuPDF), перечни олимпиад РСОШ (PDF)~\cite{rsosh}, страницы программ \texttt{ba.hse.ru}~\cite{hsebachelor} (HTML, trafilatura), FAQ (HTML, BeautifulSoup по селектору \texttt{.faq\_\_question}) и словарь абитуриента (HTML). Полный объём --- 537 чанков, распределение по типам приведено в таблице \ref{tab:kb_sources}. PDF скачиваются курированным whitelist'ом из 11 актуальных документов (\texttt{ba.hse.ru/docs}) --- sitemap не используется из-за архивных материалов прошлых циклов. FAQ переписан абитуриентским языком и даёт почти прямые матчинги на типовых запросах (см. эксп. 0). Словарь обогащается развёрткой аббревиатур (<<ИД --- индивидуальные достижения>> $\rightarrow$ пара <<Индивидуальные достижения (ИД)>>).

\subsection{Векторное хранилище}
ChromaDB~\cite{chromadb} в embedded-режиме, persistent на диске, cosine similarity, фильтры по метаданным через \texttt{where}-клаузу. Корпус $\sim 10^3$ чанков и нагрузка 50--100 запросов в день покрываются без отдельного сервиса. Коллекция --- \texttt{hse\_admission} (бакалавриат, Москва).

\subsection{Эмбеддер}
\texttt{BAAI/bge-m3}~\cite{bgem3}, 1024-dim, выбран по эксп. 3 (раздел 7.4): лучший hit@1 и MRR из четырёх кандидатов. Загружается через \texttt{sentence-transformers}~\cite{sentencetransformers} с Metal-ускорением.

\subsection{LLM-генератор}
Ollama~\cite{ollama}, локальный рантайм с Metal-ускорением на M4. Переключение моделей через \texttt{OLLAMA\_MODEL} без перезапуска бота. Системный промпт ограничивает контекстом и инструктирует честно сообщать об отсутствии информации в источниках.

\subsection{Telegram-интерфейс}
Aiogram 3.x~\cite{aiogram} в long polling. REST-API между ботом и пайплайном не введена --- прямой импорт \texttt{rag.pipeline} исключает сетевой overhead и упрощает запуск.

\subsection{Структура чанка}
JSON Lines (\texttt{data/processed/kb.jsonl}), индексация в ChromaDB как (\texttt{embedding}, \texttt{document}, \texttt{metadata}):

\begin{verbatim}
{"chunk_id": "hse_ba_admission_2026_p3_c2", "text": "...",
 "source_url": "https://ba.hse.ru/...", "source_type": "web",
 "doc_title": "Правила приёма 2026", "doc_type": "rules",
 "section": "Приём документов", "year": 2026}
\end{verbatim}

Поле \texttt{doc\_type} принимает значения \texttt{rules}, \texttt{bvi}, \texttt{appendix}, \texttt{faq}, \texttt{program}, \texttt{glossary}, \texttt{olymp\_list} и используется для фильтрации в инструменте агента (раздел 7.8).

\section{Подход к тестированию и обоснование метрик качества}

\subsection{Тестовый набор}
120 вопросов на русском с разметкой: 60 из реальных FAQ и каналов абитуриентов, 30 синтетических на основе чанков KB с ручной проверкой, 30 целенаправленно подобранных по проблемным категориям (определения, обобщённые формулировки, multi-hop). Поля: \texttt{qid}, \texttt{question}, \texttt{reference\_answer}, \texttt{relevant\_chunk\_ids}, \texttt{category}, \texttt{difficulty}, \texttt{source}, \texttt{needs\_multi\_hop}. Распределение по сложности: easy --- 49, medium --- 61, hard --- 10; multi-hop --- 10.

\subsection{LLM-as-judge}
\texttt{relevant\_chunk\_ids} размечаются автоматически: для каждого вопроса retriever возвращает top-10 кандидатов, \texttt{Qwen2.5:7b} через Ollama выбирает идентификаторы с фактическим ответом. Ручная разметка 120 вопросов потребовала бы 4--6 часов на каждую переразметку, а переразметки выполняются после каждой смены чанкера или эмбеддера.

\subsection{Стратифицированный сплит}
dev --- 85 вопросов ($\approx$70 \%) для всех экспериментов, test --- 35 ($\approx$30 \%) для финального прогона. Стратификация по \texttt{(category, difficulty)} с \texttt{seed = 42}. Test открывается ровно один раз.

\subsection{Retrieval-метрики}
\begin{itemize}
    \item Recall@k --- доля вопросов, для которых хотя бы один из \texttt{relevant\_chunk\_ids} попал в top-k. Главная метрика, целевое значение Recall@5 $\geq 0.75$.
    \item Hit@k --- индикатор <<top-k содержит хотя бы один релевантный>>. На бинарной разметке совпадает с Recall@k.
    \item MRR --- 1/rank первого релевантного, усреднённое по вопросам. Тонкая метрика для близких по recall конфигураций.
\end{itemize}

nDCG не используется: разметка бинарная, на ней nDCG вырождается в комбинацию precision@k и MRR.

\subsection{Generation-метрики}
\begin{itemize}
    \item ROUGE-L~\cite{rouge} --- лексическое перекрытие через длиннейшую общую подпоследовательность. Главный авто-сигнал.
    \item BERTScore~\cite{bertscore} (multilingual) --- семантическая близость, спасает на перефразированных ответах.
    \item Faithfulness --- доля утверждений в ответе, подкреплённых контекстом. Ручная разметка на 20--30 ответах финалистов.
    \item Answer relevance --- экспертная оценка 1--5, два прохода с интервалом, при расхождении более 1 балла --- третий проход.
\end{itemize}

BLEU не используется (плохо работает на ответах до 50 слов).

\subsection{Методологическое ограничение}
\texttt{relevant\_chunk\_ids} получены LLM-as-judge на top-10 baseline-retriever'а (\texttt{e5-large}~\cite{e5}). Это shadow ground truth: Recall@10 baseline'а равен 1.0 по построению, Recall@5 близок к 1.0. Числа показывают согласие judge с retriever'ом, а не реальное качество retrieval'а. Эксп. 0 (расширение KB) и эксп. 1, 2 (preprocessing, chunking при переразметке) валидны. Эксп. 3 (embedder) biased по абсолютному recall@5, но MRR и hit@1 информативны. Эксп. 3.5 и 4 --- валидны как сравнения относительно dense-baseline. Основной сигнал во всех экспериментах --- MRR и Recall@1, не Recall@5.

\section{Результаты тестирования и сравнение с аналогами}

Все эксперименты на \texttt{golden\_dev} (85 вопросов). Для каждой смены чанкера или эмбеддера --- полная переиндексация и переразметка \texttt{relevant\_chunk\_ids}.

\subsection{Эксперимент 0. Расширение базы знаний}

Цель --- оценить вклад каждого типа источника. Эмбеддер \texttt{e5-large}~\cite{e5}, чанкер recursive 2000/200, без реранкера.

\begin{table}[ht]
\caption{Результаты Эксп. 0 --- Расширение базы знаний.}
\label{tab:exp0}
\centering
\begin{tabular}{lccccc}
\toprule
Конфигурация & \# чанков & hit@5 & recall@5 & MRR & $\Delta$ hit@5 \\
\midrule
минимальная (rules + appendix) & 201 & 0.346 & 0.258 & 0.298 & --- \\
+ olymp + bvi & 378 & 0.383 & 0.315 & 0.321 & +0.037 \\
+ faq & 487 & 0.877 & 0.790 & 0.653 & +0.494 \\
+ glossary & 502 & 0.926 & 0.894 & 0.720 & +0.049 \\
full (+ program) & 537 & 1.000 & 1.000 & 0.759 & +0.074 \\
\bottomrule
\end{tabular}
\end{table}

FAQ даёт основной прирост (+0.494 hit@5) --- 109 готовых пар Q--A покрывают обобщённые формулировки. Словарь даёт лучшее соотношение <<ценность/объём>> (+0.049 за 15 чанков). Дальнейшие эксперименты --- на полном корпусе.

\subsection{Эксперимент 1. Предобработка}

Цель --- оценить вклад дедупликации и heading-разделителей.

\begin{table}[ht]
\caption{Результаты Эксп. 1 --- Предобработка.}
\label{tab:exp1}
\centering
\begin{tabular}{lccccc}
\toprule
Конфигурация & \# чанков & hit@5 & recall@5 & MRR & recall@1 \\
\midrule
basic & 576 & 0.877 & 0.798 & 0.640 & 0.309 \\
+ dedup & 537 & 1.000 & 1.000 & 0.759 & 0.365 \\
+ headings (default) & 537 & 1.000 & 1.000 & 0.759 & 0.365 \\
\bottomrule
\end{tabular}
\end{table}

Дедупликация критична ($-0.122$ hit@5 без неё). Heading-разделители на данном корпусе нейтральны: PDF почти не содержат markdown-заголовков. Оставлены как дефолт --- не вредят, могут пригодиться при добавлении структурированных web-источников.

\subsection{Эксперимент 2. Чанкинг}

Цель --- выбор метода и размера чанка.

\begin{table}[ht]
\caption{Результаты Эксп. 2 --- Чанкинг.}
\label{tab:exp2}
\centering
\begin{tabular}{lcccc}
\toprule
Метод & \# чанков & hit@1 & hit@5 & MRR \\
\midrule
fixed 512 chars & 537 & 0.600 & 0.988 & 0.754 \\
fixed 1024 chars & 325 & 0.563 & 1.000 & 0.729 \\
recursive (default) & 537 & 0.613 & 1.000 & 0.765 \\
\bottomrule
\end{tabular}
\end{table}

Recursive splitter с heading-разделителями выигрывает по MRR (0.765). Большие чанки (1024) замывают релевантность шумом ($-0.036$ MRR). Финал --- recursive, 2000 символов, overlap 200.

\subsection{Эксперимент 3. Сравнение эмбеддеров}

Цель --- выбор multilingual-эмбеддера для русского admission-домена.

\begin{table}[ht]
\caption{Результаты Эксп. 3 --- Embedder.}
\label{tab:exp3}
\centering
\begin{tabular}{lccccc}
\toprule
Embedder & Dim & hit@1 & hit@5 & MRR & n\_scored \\
\midrule
e5-base & 768 & 0.679 & 0.988 & 0.795 & 81 \\
e5-large & 1024 & 0.600 & 0.988 & 0.754 & 80 \\
LaBSE-en-ru & 768 & 0.554 & 1.000 & 0.709 & 65 \\
bge-m3 & 1024 & 0.700 & 1.000 & 0.807 & 80 \\
\bottomrule
\end{tabular}
\end{table}

\texttt{bge-m3}~\cite{bgem3} побеждает по hit@1 и MRR --- современный SOTA multilingual-эмбеддер с нативной поддержкой sparse+dense. Архитектура и тренировочные данные важнее размерности: \texttt{e5-base}~\cite{e5} (768) сильнее \texttt{e5-large} (1024) на малом корпусе. \texttt{LaBSE-en-ru} сильно проседает несмотря на специализацию. Все последующие эксперименты --- на \texttt{bge-m3}.

\subsection{Эксперимент 3.5. Гибридный поиск}

Цель --- оценить прирост от BM25~\cite{bm25} и LLM-перефразирования поверх dense.

\begin{table}[ht]
\caption{Результаты Эксп. 3.5 --- Hybrid retrieval и query rewriting.}
\label{tab:exp35}
\centering
\begin{tabular}{lcccc}
\toprule
Метод & hit@1 & hit@5 & recall@5 & MRR \\
\midrule
dense (bge-m3) & 0.638 & 1.000 & 1.000 & 0.775 \\
BM25 only & 0.188 & 0.425 & 0.356 & 0.290 \\
hybrid RRF & 0.438 & 0.750 & 0.683 & 0.551 \\
hybrid + rewrite & 0.388 & 0.700 & 0.617 & 0.518 \\
\bottomrule
\end{tabular}
\end{table}

BM25~\cite{bm25} слаб на маленьком русском корпусе без морфологической нормализации. Hybrid RRF (k=60) проседает на 0.224 MRR: слабые BM25-кандидаты вытесняют сильные dense~\cite{dpr}. Rewrite добавляет шум и +3 сек latency. Значительная часть просадки --- methodological bias раздела 6.6 (judge выбирал из dense top-5). Решение --- чистый dense; гибрид и rewrite оставлены в коде за переменной \texttt{RETRIEVE\_METHOD}.

\subsection{Эксперимент 4. Реранкер}

Цель --- оценить cross-encoder~\cite{sbert} поверх dense top-20.

\begin{table}[ht]
\caption{Результаты Эксп. 4 --- Реранкер.}
\label{tab:exp4}
\centering
\begin{tabular}{lccccc}
\toprule
Конфигурация & hit@1 & hit@5 & recall@5 & MRR & Latency, сек/запрос \\
\midrule
none (baseline) & 0.638 & 1.000 & 1.000 & 0.775 & около 0.2 \\
bge-reranker-v2-m3 & --- & --- & --- & --- & $> 15$ (timeout) \\
mmarco-mMiniLMv2 & не прогнано & --- & --- & --- & --- \\
\bottomrule
\end{tabular}
\end{table}

\texttt{sentence\_transformers.CrossEncoder.predict()} без batching и MPS на M4 CPU даёт 2--4 сек на пару --- 40--80 сек на 20 пар при бюджете 10 сек. Реранкер на CPU непригоден; перепроверка на GPU --- раздел 8.2.

\subsection{Эксперимент 5. Генерация}

Coordinate descent по трём осям (модель, промпт, температура) --- 12 конфигураций вместо полного grid (36). Retrieved top-k для каждого dev-вопроса кэшируется в \texttt{data/processed/retrieved\_cache.jsonl}.

Подэтап 5A --- выбор модели при simple prompt, T = 0.3.

\begin{table}[ht]
\caption{Результаты эксп. 5A --- выбор LLM-генератора (simple prompt, T = 0.3).}
\label{tab:exp5a}
\centering
\begin{tabular}{lccc}
\toprule
Модель & ROUGE-L & BERTScore F1 & avg latency, сек \\
\midrule
qwen2.5:7b-instruct & 0.062 & 0.682 & 2.13 \\
gemma4:e4b & 0.012 & --- & 3.88 \\
mistral:latest & 0.070 & 0.707 & 1.69 \\
llama3:latest & 0.071 & 0.704 & 1.16 \\
\bottomrule
\end{tabular}
\end{table}

BERTScore F1~\cite{bertscore} для \texttt{gemma4:e4b} не рассчитана из-за конфликта transformers и bert\_score. Модель в большинстве ответов воспроизводила <<В моих источниках этого нет>> из системной инструкции без попытки извлечь ответ --- избыточная консервативность 4B-модели на справочных запросах. Топ-2 (\texttt{mistral:latest}, \texttt{llama3:latest}) идут в 5B.

Подэтап 5B --- выбор промпта (simple, citing, few-shot) на топ-2 моделях, T = 0.3.

\begin{table}[ht]
\caption{Результаты эксп. 5B --- выбор промпт-стратегии (T = 0.3).}
\label{tab:exp5b}
\centering
\begin{tabular}{llcc}
\toprule
Модель & Промпт & ROUGE-L & BERTScore F1 \\
\midrule
mistral:latest & simple & 0.084 & 0.710 \\
mistral:latest & citing & 0.059 & 0.683 \\
mistral:latest & fewshot & 0.068 & 0.707 \\
llama3:latest & simple & 0.068 & 0.706 \\
llama3:latest & citing & 0.047 & 0.673 \\
llama3:latest & fewshot & 0.070 & 0.692 \\
\bottomrule
\end{tabular}
\end{table}

Победитель --- \texttt{mistral:latest} + simple. Citing снижает BERTScore: модель расходует ёмкость на форматирование \texttt{[N]} в ущерб семантической полноте.

Подэтап 5C --- температура на \texttt{mistral:latest} + simple.

\begin{table}[ht]
\caption{Результаты эксп. 5C --- подбор температуры (mistral:latest + simple).}
\label{tab:exp5c}
\centering
\begin{tabular}{lccc}
\toprule
Temperature & ROUGE-L & BERTScore F1 & avg latency, сек \\
\midrule
0.0 & 0.075 & 0.7137 & 1.53 \\
0.3 & 0.076 & 0.7121 & 1.60 \\
0.7 & 0.061 & 0.7078 & 1.61 \\
\bottomrule
\end{tabular}
\end{table}

Победитель --- T = 0.0. На справочных запросах детерминированное декодирование не вносит случайных словоформ, латентность ниже.

\subsection{AgenticRAG}

Гипотеза --- multi-hop вопросы выиграют от ReAct-петли~\cite{react} с инструментом \texttt{search\_kb(query, doc\_type=None, k=5)}. Реализация --- тонкий агент на Python без фреймворков, лимит 3 шага. Метрики --- на подмножестве \texttt{needs\_multi\_hop = true} (7 вопросов).

\begin{table}[ht]
\caption{Результаты эксп. AgenticRAG на подмножестве \texttt{needs\_multi\_hop = true} (7 вопросов).}
\label{tab:exp_agentic}
\centering
\begin{tabular}{lcc}
\toprule
Метрика & Baseline (single-shot RAG) & AgenticRAG \\
\midrule
avg ROUGE-L & 0.174 & 0.253 \\
avg latency, сек & 2.98 & 2.34 \\
avg steps & 1 (по построению) & 1.0 \\
\bottomrule
\end{tabular}
\end{table}

\texttt{avg\_steps = 1.0} означает, что \texttt{mistral:latest} ни разу не вызвал инструмент и сразу выдавал \texttt{Final Answer} --- фактически работал как vanilla-LLM. Прирост ROUGE-L (+0.079) и снижение latency ($-0.64$ сек) объясняются меньшим шумом контекста: single-shot RAG передаёт пять чанков, часть нерелевантна для multi-hop, vanilla-LLM даёт более усреднённый, но семантически близкий ответ.

Гипотеза о выигрыше от ReAct~\cite{react} не подтверждена для моделей 7--8B с локальным инференсом --- они слабо следуют tool-инструкциям. Негативный результат указывает на проблему шума в retrieved-контексте на сложных вопросах. Направления продолжения --- cross-encoder реранкер с агрессивным top-1, multi-query retrieval, RAG-fusion, либо более крупная модель ($\geq 14$B) с дообучением на ReAct-датасете.

\subsection{Итоговая конфигурация}
\begin{itemize}
    \item Корпус --- 537 чанков, шесть типов источников.
    \item Препроцессинг --- Unicode-нормализация, dedup, heading-разделители.
    \item Чанкер --- RecursiveCharacterTextSplitter, 2000 символов, overlap 200.
    \item Эмбеддер --- \texttt{BAAI/bge-m3}, 1024-dim.
    \item Хранилище --- ChromaDB, cosine, persistent.
    \item Retrieval --- чистый dense, top-5; реранкер отсутствует.
    \item Генератор --- \texttt{mistral:latest} через Ollama, simple prompt, T = 0.0.
\end{itemize}

Метрики на \texttt{golden\_dev}: ROUGE-L = 0.075, BERTScore F1 = 0.714, avg latency 1.53 сек, p95 = 3.44 сек.

\subsection{Сравнение с аналогами}

Сопоставление с двумя типовыми решениями: статический FAQ \texttt{ba.hse.ru/faq}~\cite{hsebachelor} (109 пар Q--A, ручной поиск глазами, не агрегирует данные из нескольких источников) и универсальные LLM-чат-боты без RAG (ChatGPT~\cite{gpt}, YandexGPT --- хорошо понимают переформулировки, но галлюцинируют нормативные детали и не различают актуальные правила и архивные).

\begin{table}[ht]
\caption{Сравнение разработанной системы с двумя типовыми аналогами.}
\label{tab:comparison}
\centering
\begin{tabular}{lcccc}
\toprule
Решение & hit@5 на golden\_dev & Переформулировки & Цитирование & Hallucinations \\
\midrule
FAQ ВШЭ (статический) & 0.346 (rules+appendix) & низкое & да & нет \\
Универсальный LLM без RAG & не измерено & высокое & нет & высокий риск \\
Разработанная система & 1.000 & высокое & да & низкий \\
\bottomrule
\end{tabular}
\end{table}

Разработанная система использует FAQ как один из источников и одновременно покрывает ограничения статического интерфейса.

\section{Заключение}

\subsection{Основные результаты работы}
Финальная конфигурация (\texttt{bge-m3}~\cite{bgem3} + dense top-5 + \texttt{mistral:latest} + simple prompt + T = 0.0) показывает hit@5 = 1.0 и MRR = 0.807 на dev, ROUGE-L~\cite{rouge} = 0.075, BERTScore F1~\cite{bertscore} = 0.714. Главный архитектурный вклад --- шестикомпонентный корпус 537 чанков; технический --- выбор \texttt{bge-m3} (эксп. 3) и \texttt{mistral:latest} (эксп. 5).

Сопоставление dev и test.

\begin{table}[ht]
\caption{Сравнение метрик финальной конфигурации на dev и test.}
\label{tab:dev_test}
\centering
\begin{tabular}{lccc}
\toprule
Метрика & dev (85 вопросов) & test (35 вопросов) & $\Delta$ \\
\midrule
ROUGE-L & 0.075 & 0.024 & $-0.051$ \\
BERTScore F1 & 0.714 & 0.697 & $-0.017$ \\
avg latency, сек & 1.53 & 1.67 & +0.14 \\
p95 latency, сек & 3.44 & 3.40 & стабильно \\
\bottomrule
\end{tabular}
\end{table}

BERTScore F1 на test стабилен ($-2.4$ \%). ROUGE-L просел на 67 \%: метрика лексически чувствительна, штрафует за синонимы и порядок слов (для русского без морфологического стеммера особенно строго), а эталоны dev и test писались независимо --- все эксперименты косвенно подгоняли пайплайн под лексику dev-эталонов. Стабильность BERTScore F1 более информативна.

Достижение целей.

\begin{table}[ht]
\caption{Сопоставление целевых и фактических значений на test.}
\label{tab:goals}
\centering
\begin{tabular}{lccc}
\toprule
Цель & Целевое значение & Факт на test & Достигнута \\
\midrule
Recall@5 (на golden dev) & $\geq 0.75$ & 1.000 на dev & формально да \\
ROUGE-L по ответам & $\geq 0.35$ & 0.024 & нет \\
BERTScore F1 & не задано & 0.697 & средняя для RU \\
Latency p95 & $\leq 10$ сек & 3.40 сек & да \\
Hallucination rate & $\leq 10$ \% & требует разметки & отложено \\
Воспроизводимость & README + одна команда & да & да \\
\bottomrule
\end{tabular}
\end{table}

Цели по retrieval, latency и воспроизводимости достигнуты. ROUGE-L не достигнута; BERTScore F1, как более подходящая для русского, показала разумный результат. Faithfulness и expert relevance отложены в раздел 8.2. Методологически фиксируется ограничение shadow ground truth --- основной сигнал во всех сравнениях --- MRR и hit@1.

\subsection{Перспективы дальнейшей деятельности}
\begin{itemize}
    \item Честный ground truth --- независимая разметка 30 вопросов через альтернативный pool top-50 от трёх retrieval-методов, ручная faithfulness (0/1) и expert relevance (1--5) на финальной конфигурации.
    \item Retrieval --- cross-encoder реранкер на GPU (T4/A100), sparse-режим \texttt{bge-m3} вместо отдельного BM25~\cite{bm25}, contrastive fine-tune эмбеддера на парах \texttt{question}--\texttt{relevant\_chunk}.
    \item Снижение шума контекста на multi-hop --- multi-query retrieval, RAG-fusion, агрессивный top-1 после реранкера.
    \item AgenticRAG на моделях $\geq 14$B с дообучением на ReAct-датасете~\cite{react} или расширением инструментов (\texttt{compare\_programs}, \texttt{lookup\_olymp\_by\_subject}).
    \item Multi-turn диалог с памятью, feedback-loop через оценки в Telegram, расширение скоупа --- магистратура, кампусы СПб/НН/Пермь, multi-modal retrieval (таблицы и схемы из PDF).
\end{itemize}

\newpage
\printbibliography[heading=bibintoc]

\newpage
\appendix
\section{Терминологический словарь}

Определения по словарю абитуриента ВШЭ (\texttt{ba.hse.ru/vocabulary})~\cite{hsebachelor} и Правилам приёма~\cite{hsebachelor}.

\textbf{БВИ} (без вступительных испытаний) --- право на зачисление без сдачи общих ВИ; даётся победителям и призёрам ВсОШ по профильному предмету и призёрам перечневых олимпиад первого--третьего уровня при максимальном балле по профильному ЕГЭ.

\textbf{ВсОШ} --- Всероссийская олимпиада школьников; победа в заключительном этапе даёт БВИ по профильному предмету на программы РФ.

\textbf{РСОШ}~\cite{rsosh} --- Российский совет олимпиад школьников; ежегодно утверждает перечень олимпиад, дающих льготы. Перечень делится на три уровня.

\textbf{ИД} --- индивидуальные достижения; до 10 дополнительных баллов за олимпиады непрофильного перечня, аттестат с отличием, итоговое сочинение, ГТО.

\textbf{ОРМ} --- особое право мобильности; льгота для отдельных категорий (дети-сироты, инвалиды), бюджет в пределах квоты при минимальных баллах.

\textbf{ДВИ} --- дополнительные ВИ, проводимые вузом поверх ЕГЭ (на ВШЭ --- для <<Дизайна>>, <<Журналистики>>, <<Восточной и африканской филологии>>).

\textbf{КЦП} --- контрольные цифры приёма; число бюджетных мест на программу, делится на общие, целевую, особую и отдельную квоты.

\textbf{Целевой приём} --- поступление по договору с организацией-заказчиком, оплачивающей обучение; конкурс --- внутри квоты.

\textbf{Конкурсная группа} --- совокупность программ, ранжируемых единым списком. На ФКН --- отдельные программы, на ряде факультетов одно направление объединяет несколько программ.

\textbf{Зелёная волна} --- приоритетное зачисление БВИ-абитуриентов и победителей олимпиад уровня вуза до основных этапов.

\textbf{Лист ожидания} --- список не прошедших в основной волне, имеющих право на зачисление при освобождении мест.

\textbf{Ранжирование} --- упорядочение по убыванию суммы ЕГЭ + ДВИ + ИД; тай-брейк: профильный ЕГЭ, затем русский, затем дата согласия.

\textbf{Скидки} --- снижение стоимости коммерческого обучения по шкале 25 / 50 / 70 \% в зависимости от баллов.

\textbf{Проходной балл} --- балл последнего зачисленного на бюджет в прошлом цикле; ориентир, не гарантия.

\textbf{Минимальный балл} --- нижняя граница приёма ЕГЭ/ДВИ; ежегодно устанавливается вузом.

\textbf{ЦТ} --- централизованное тестирование (Беларусь); засчитывается ВШЭ для граждан Беларуси.

\section{Список сокращений}

\textbf{API} --- Application Programming Interface. \textbf{BERTScore} --- метрика семантической близости на BERT-эмбеддингах. \textbf{БВИ} --- без вступительных испытаний. \textbf{BM25} --- Best Match 25, ranking-функция лексического поиска. \textbf{ВИ} --- вступительные испытания. \textbf{ВсОШ} --- Всероссийская олимпиада школьников. \textbf{ВШЭ} --- Высшая школа экономики. \textbf{GPU} --- Graphics Processing Unit. \textbf{HTML} --- HyperText Markup Language. \textbf{ДВИ} --- дополнительные ВИ. \textbf{ЕГЭ} --- Единый государственный экзамен. \textbf{ИД} --- индивидуальные достижения. \textbf{JSONL} --- JSON Lines. \textbf{KB} --- Knowledge Base. \textbf{КЦП} --- контрольные цифры приёма. \textbf{LLM} --- Large Language Model. \textbf{MRR} --- Mean Reciprocal Rank. \textbf{ОРМ} --- особое право мобильности. \textbf{PDF} --- Portable Document Format. \textbf{RAG} --- Retrieval-Augmented Generation. \textbf{ROUGE-L} --- Recall-Oriented Understudy for Gisting Evaluation (LCS-вариант). \textbf{RRF} --- Reciprocal Rank Fusion. \textbf{РСОШ} --- Российский совет олимпиад школьников. \textbf{СПО} --- среднее профессиональное образование. \textbf{ЦТ} --- централизованное тестирование. \textbf{CPU} --- Central Processing Unit.

\section{Описание исходных данных}

База знаний собрана скриптом \texttt{scripts/ingest\_all.py} с whitelist'ом 11 PDF и фиксированным списком HTML-URL.

\begin{table}[ht]
\caption{Источники KB.}
\label{tab:kb_sources}
\centering
\begin{tabular}{llll}
\toprule
Источник & doc\_type & URL & Объём в KB \\
\midrule
Правила приёма (PDF) & rules & ba.hse.ru/docs & 90 чанков \\
Приложения к правилам (PDF) & appendix & ba.hse.ru/docs & 111 чанков \\
Перечни олимпиад (PDF) & olymp\_list & ba.hse.ru/docs, rsr-olymp.ru & 177 чанков \\
FAQ ВШЭ (HTML) & faq & ba.hse.ru/faq & 109 чанков \\
Словарь абитуриента (HTML) & glossary & ba.hse.ru/vocabulary & 15 чанков \\
Страницы программ (HTML) & program & ba.hse.ru/* & 35 чанков \\
Всего &  &  & 537 чанков \\
\bottomrule
\end{tabular}
\end{table}

Тестовый набор 120 вопросов --- \texttt{data/golden\_set.jsonl}. Стратифицированный split 85/35 (\texttt{seed = 42}) сохранён в \texttt{data/golden\_dev.jsonl} и \texttt{data/golden\_test.jsonl}. Test открывается ровно один раз.

\end{document}
